\section{An isolated condition must survive internal graph modification}
\label{exec13:isolated-condition}

Preserving the node set at the adapter boundary is necessary but does not
establish that every internal graph operation preserves it. The inspected y0
revision's \texttt{remove\_in\_edges} rebuilds a graph from the intervention
vertices and the endpoints of surviving edges. It does not explicitly seed
every original vertex. An isolated vertex outside the intervention can
therefore disappear. The ID recursion later requests ancestors of its outcome
set in this modified graph. These are source-level observations about the
same pinned revision, not results from an installed package trial
\citep{exec13_y0_graph,exec13_y0_id}.

Consider the supplied DAG with \(A\to Y\) and an isolated condition \(Z\).
The requested object is \(p(y\mid\operatorname{do}(a),z)\). The proposed
joint route begins with \(p(y,z\mid\operatorname{do}(a))\), retaining
\(Y,Z\) as its outcomes. All three vertices are ancestors of this outcome
set or members of it before graph modification, so the initial ancestor
restriction does not remove \(Z\). In the subsequent source path, removing
incoming edges to \(A\) reconstructs vertices \(A,Y\) from the intervention
and the edge \(A\to Y\). The ancestor request for \(Z\) then concerns a
vertex that is absent. An actual installed-engine run is required to record
its precise exception and response; the source trace already establishes
why checking only the initial conversion is insufficient
\citep{exec13_isolated_reference}.

This failure mechanism is distinct from a failure of causal identification.
For an exact reference, let \(A,Z\) be independent fair binary variables and
let an independent outcome mechanism satisfy
\[
 P(Y=1\mid\operatorname{do}(A=a))=\tfrac14+\tfrac12a.
\]
The eight observed cells have positive probability. Under the supplied
causal DAG and its independent mechanisms,
\[
 p(y,z\mid\operatorname{do}(a))=p(y\mid a)p(z),\qquad
 P(Y=1\mid\operatorname{do}(a),z)=\tfrac14+\tfrac12a
\]
on both supported values of \(z\). The effect is \(1/2\) within either
condition. Independent rational enumeration reproduces this result. These
are stipulated mechanisms, not inferred bank facts or fitted package
outputs \citep{exec13_isolated_reference}.

The admission decision is consequently more specific. The proposed joint
route needs tests for isolated conditions and outcomes at every internal
graph transformation, in addition to the front-door, bow and collider
references. A reviewed patch that retains all vertices would need an actual
engine comparison and regression checks before admission. An adapter may
instead implement an independently justified isolated-variable factorization
for the exact supported graph class, but must not discard a condition merely
because it appears inconvenient. A condition with an incident edge does not
qualify for the isolated-variable shortcut. Until a route passes these
checks, a missing-node exception is an unavailable computation with the
original question retained. It is neither a proof of nonidentification nor
a successful empty result. The public proposal can therefore recommend the
restricted integration while keeping its unresolved implementation obligation
visible.
\label{exec13:isolated-condition-end}
